\headline{{\textbf{\Large\sffamily Seasonal Care Tips:}}\\  
          {\textbf{\large\sffamily Mid\--July to Mid\--August}}}
\label{2025-07-care}
\noindent  
From mid\--June to mid\--July, bonsai in our area face the full force of summer's intensity.
With prolonged daylight and soaring temperatures---exacerbated by recent heatwaves---trees are in active growth but also under significant stress.
Balancing hydration, pest control, and sun protection becomes paramount.
Pay particular attention to microclimates in your garden; shaded spots or sheltered corners can offer crucial respite during extreme heat.
Here's how to navigate these demanding weeks while keeping your trees healthy and thriving.

\medskip  
\topic{Growth Management in High Heat}  
\noindent  
Deciduous species like maples and hornbeams will have lush foliage, but avoid heavy pruning during heatwaves.
Instead, thin out overcrowded leaves to improve airflow and reduce moisture loss.
For pines, the second flush from earlier decandling should emerge---trim new shoots lightly to maintain shape, but delay aggressive work until temperatures cool.
Consider postponing any non\--essential pruning until evenings when temperatures drop, reducing stress on the tree.

Junipers and conifers benefit from pinching back elongated growth, but refrain from cutting into old wood, as recovery slows in extreme heat.
Flowering species, such as wisteria or chaenomeles, should be deadheaded to conserve energy.
Tropical trees (ficus, bougainvillea) will grow vigorously; guide their structure with light pruning and wiring.

\medskip  
\topic{Watering: Combatting Drought Stress}  
\noindent  
Heatwaves demand vigilance.
Check soil moisture \textbf{twice daily}---small pots may need watering by mid\--morning and again in the evening.
Water thoroughly until it drains freely, ensuring the entire root mass is saturated.
For trees in fast\--draining mixes (e.g., pines or junipers), consider a morning soak in a shallow tray to boost hydration.

Grouping trees together or placing them on damp gravel trays can raise humidity.
If away briefly, use capillary matting or self\--watering spikes, but test first to ensure even moisture. \textbf{Never water scorching\--hot soil}---cool pots slightly with shade or a misting spray first to avoid root shock.
In prolonged dry spells, a layer of fine akadama or sphagnum moss atop the soil can help retain moisture without suffocating roots.

\medskip  
\topic{Fertilising with Caution}  
\noindent  
Reduce nitrogen\--heavy feeds for deciduous trees to slow soft growth, which is vulnerable to heat stress.
Conifers still benefit from balanced, slow\--release fertiliser, but dilute liquid feeds to half\--strength to prevent root burn.
For flowering/fruiting species, switch to a potassium\--rich formula to support next year's buds.

Organic options (e.g., seaweed extract) can bolster resilience but apply sparingly---excess may attract pests in the heat.
Always water the soil before fertilising to protect roots.
If using granular fertiliser, place it slightly farther from the trunk than usual to avoid concentrated salts in the root zone during dry spells.

\medskip
\topic{Pests and Diseases in the Heat}  
\noindent  
Spider mites, aphids, and scale thrive in dry conditions.
Inspect foliage undersides for stippling or webbing.
Blast pests off with water or treat with insecticidal soap (apply at dawn/dusk to avoid leaf scorch).
For severe infestations, neem oil can be effective but test on a small area first---some species (e.g., azaleas) are sensitive.

Fungal issues like powdery mildew may arise in humid spells.
Prune dense foliage to improve airflow and remove affected leaves promptly.
Avoid overhead watering in the evening to keep foliage dry overnight.
Consider preventative treatments for vulnerable species, applying fungicides in the cooler early morning hours for best absorption.

\medskip  
\topic{Shielding Trees from Scorching Sun}  
\noindent  
Delicate species (Japanese maples, young saplings) may need shade cloth or dappled sunlight during peak afternoon heat.
Even sun\--loving pines can suffer---watch for needle browning.
Light\--coloured pots or stands help reflect heat, and elevating pots improves airflow.
Rotate trees periodically to ensure all sides receive equal light and avoid lopsided growth.

Tropicals relish warmth but shelter them from drying winds.
For bonsai on patios or stone surfaces, place them on wooden boards to insulate roots from radiant heat.
During heatwaves, temporary shading with bamboo screens or horticultural fleece can prevent leaf scorch without drastically reducing light levels.

\medskip  
\topic{Wiring and Structural Tweaks}  
\noindent  
Rapid growth can cause wires to bite within weeks.
Check wired branches every 5–7 days and loosen as needed.
Remove wires from deciduous trees if scars appear; conifers may tolerate them longer but monitor closely.
Avoid major styling---save heavy bending or repotting for autumn.
If adjusting branches, work in the early morning when tissues are more pliable and less prone to cracking.

\medskip  
\topic{Early Preparations for Late Summer}  
\noindent  
By mid\--July, start planning ahead:  
\begin{itemize}
    \item \textbf{Reduce fertilisation} gradually for deciduous trees to harden off new growth.  

    \vspace{-0.5em}
    \item \textbf{Collect seeds} (e.g., acers, elms) for autumn sowing, storing them in cool, dry conditions.  

    \vspace{-0.5em}
    \item \textbf{Rehydrate or replace dried moss} to retain soil moisture and aesthetics.
    Begin scouting for autumn repotting candidates, noting which trees might benefit from root work as the season transitions.
\end{itemize}

\vspace{-0.5em}
\topic{Embracing the Season's Challenges}
\noindent  
Despite the hurdles, summer offers moments of joy---vibrant foliage, the hum of bees around flowering trees, and the satisfaction of nurturing resilience.
Share your heatwave strategies at the next club meeting, whether it's a shade hack or a pest remedy.
Document any surprising successes, like species that thrived despite the odds---these observations enrich our collective knowledge.
With attentive care, your bonsai will weather the heat and emerge stronger.
\closearticle
\columnbreak  
